\scanpage{53}
\mspage{44}
% archive sheet number: 53

\deleted{par c) $\exists j'\leq j$ avec $x\in X_{j'}^{\circ\prime}$, a fortiori $x\in X_{j'}^\circ$. Donc $X_i^\circ\cap X_{j'}^\circ\ne\varnothing$, donc par $b'')$, on a $i=j'$, d'où $i\leq j$, qfd.}

\uline{Exemple} Voici un exemple où on a $1'$ et $2'$, et en plus tous les $X_i^\circ\ne\varnothing$, sans avoir $3'$. \inserted{(indexation respectant la relation d'inclusion $i\leq j\Longleftrightarrow X_i\subset X_j$)}

Soit $\deleted{X=}A_0\supset A_1\supset\cdots\supset A_n\supset\cdots$ une suite \uline{strictement} \uline{décroissante} d'ens. de $X$, avec
\[
\deleted{A_\infty=}\inserted{N=}\bigcap A_n\ne\varnothing.
\]
\inserted{Soit $\deleted{\unsure{B}}A_\infty\subset X$ avec $A_\infty\cap A_0=N$, $A_\infty\not\subset A_0$ i.e. $N\ne A_\infty$.}
On prend $\deleted{\unsure{B\cap A_n}}\unsure{I|\Phi}\in\unsure{\mathcal{P}|\mathrm{P}}(X)$ défini comme ens. des $A_i$ ($i\in\mathbb{N}\cup\{\infty\}$). \deleted{\unsure{I est}} \inserted{$I\setminus\{\infty\}$} un ens \deleted{\unsure{tot.}} ordonné par inclusion, mais $\infty$ est un ``pt isolé''. On a
\[
\begin{aligned}
A_i^\circ&=A_i\setminus A_{i+1}\ne\varnothing&&\text{si }i\ne\infty\\
A_\infty^\circ&=A_\infty\ne\varnothing
\end{aligned}
\]
\deleted{\unsure{Les $A_i$}} Les $A_i^\circ$ ($i\in I$) sont mutuellement disjoints non vides, leur réunion est $A_0\inserted{\cup A_\infty}$ $(=|\unsure{\Phi|P}|)$. Mais Un $A_i$ ($i\ne\infty$) n'est pas réunion des $A_j^\circ$ avec $A_j\subset A_i$ (il manque $N$ !), il n'est pas saturé par la relation d'équivalence définie par les $A_j^\circ$ ($j\in I$).

\deleted{D'autre part,} \inserted{Un} $x\in N$ n'est pas contenu dans un plus petit élément de $\Phi$. On a
\[
\underline{A_\infty^\circ\cap A_i}\inserted{=N}\ne\varnothing,
\]
\margin{$\unsure{A_\infty^\circ\cap A_i}$}
sans pour autant avoir $A_\infty^\circ\subset A_i$ \deleted{(mais bien $A_\infty^\circ\subset A_i$}

Ici la condition \unsure{-bis $\beta$)} n'est pas satisfaite (saturation des $A_i$), alors que $\alpha$, $\gamma$) le sont.

Question : si on suppose
\[
\left\{\begin{aligned}
&\forall i\quad X_i^\circ\ne\varnothing\\
&\forall i\ne j\quad X_i^\circ\cap X_j^\circ=\varnothing\\
&\deleted{\bigcup X_i^\circ=|\Phi|}
\end{aligned}\right.
\qquad+\ \begin{gathered}\text{la saturation}\\\text{de chaque }X_i\end{gathered}
\]
\[
X_i=\bigcup_{j\in I,\ X_j^\circ\subset X_i}X_j
\]
